% TOP_PROBLEM: {"id": 3219, "title": "The three 4-flows conjecture", "category_id": 3, "category": {"id": 3, "name": "graph_theory", "display_name": "Graph Theory", "description": "Problems involving graphs, networks, and their properties.", "slug": "graph-theory", "order_index": 3, "created_at": "2026-07-31T15:26:25.671Z"}, "status": "open", "problem_number": "OPG-132", "set_id": 12}
% FIRST_POSTED: 2026-10-01
% ATTEMPT_STATUS: unresolved
% MODEL: GPT 6 Astra Ultra
% UnsolvedMath ID 3219; source catalogue code OPG-132
% !TeX program = lualatex
% Research attempt dated 2026-10-01; canonical statement preserved.
\documentclass[11pt,a4paper]{article}
\usepackage[margin=25mm]{geometry}
\usepackage{fontspec}
\setmainfont{lmroman10-regular.otf}[BoldFont=lmroman10-bold.otf,
 ItalicFont=lmroman10-italic.otf,BoldItalicFont=lmroman10-bolditalic.otf]
\usepackage{amsmath,amssymb,mathtools}
\usepackage{unicode-math}
\setmathfont{latinmodern-math.otf}
\AtBeginDocument{\let\setminus\smallsetminus}
\usepackage{xurl,hyperref}
\hypersetup{colorlinks=true,urlcolor=blue}
\setcounter{secnumdepth}{0}
\setlength{\parindent}{0pt}
\setlength{\parskip}{5pt}
\setlength{\emergencystretch}{3em}
\begin{document}
\section*{The three 4-flows conjecture}\label{3219}
{\small UnsolvedMath ID 3219 · OPG-132 · Graph Theory · OpenGarden}
\subsection{Definitions and mathematical statement}
Let $G=(V,E)$ be a finite loopless
undirected multigraph without a
bridge. A bridge is an edge
whose removal increases the
number of connected components.
For an edge set $A\subseteq E$,
write $G\setminus A=(V,E\setminus A)$;
all vertices, including any newly
isolated ones, are retained.

A nowhere-zero integer $4$-flow
on a graph $H$ consists of an
orientation and a function
$\phi:E(H)\to\{-3,-2,-1,1,2,3\}$
such that the sum entering
each vertex equals the sum
leaving it. The equations
are integer equalities, and
existence is unaffected by
reversing edges and negating
their values. On an edgeless
graph the empty function is
allowed.

The conjecture asks for an
edge partition
\[
 E=A_1\mathbin{\dot\cup}A_2\mathbin{\dot\cup}A_3
\]
such that each of the three
spanning graphs $G\setminus A_i$
admits a nowhere-zero $4$-flow.
Parts of the partition may
be empty. The three flows
are chosen separately and
need not give equal values
or directions to edges present
in more than one graph.

Every original edge appears
in exactly two of the three
flow-supporting subgraphs, because
it is deleted in exactly
one part. Equivalently, one
seeks three bounded integer
flows on $G$, allowing zero
values, such that each edge
is zero in precisely one
flow. A partition into three
subgraphs each supporting a
flow would be a different
requirement: here flows live
on the complements of parts.
\subsection{Short English statement}
Can every bridgeless graph have its edges split into three parts so deleting any one part leaves a graph with a nowhere-zero 4-flow?
\subsection{Research attempt}
\paragraph{Scope and method.}
This research attempt by GPT-6 Astra Ultra is dated 1 October 2026.
It preserves the catalogue statement and distinguishes elementary proofs
below from cited prior work. No novelty claim or formal proof-assistant
verification is made. The final paragraph states the precise remaining scope.
\paragraph{Source and the attempted strengthening.}
The source [S2], posted by DeVos, asks for a partition whose three
complements admit four-flows. We investigate the stronger demand that
each complement be Eulerian, giving a unit integer flow. This succeeds
on three-edge-colourable cubic graphs but fails even on the Petersen
graph. That failure explains a limitation of the proposed method, not
a negative answer to the actual four-flow question.

\paragraph{Immediate and constructive positive cases.}
If $G$ already has a nowhere-zero four-flow, put $A_1=E(G)$ and
$A_2=A_3=\varnothing$. The complements are the empty graph, $G$, and $G$.
The empty flow is allowed on the empty edge set, so this is valid under
the catalogue's explicit convention permitting empty parts.
In particular every Eulerian graph works by orienting Euler tours and
assigning unit values. This first construction need not make all three
parts nonempty; imposing that extra condition would alter the question.

There is a more symmetric construction for a cubic graph with a proper
three-edge-colouring. Let $A_i$ be its colour classes. At every vertex,
deleting any one colour leaves degree two. Each complement is therefore
a disjoint union of cycles, which can be oriented cyclically with value
one on every edge. Each edge occurs in exactly two complements as the
question requires, and all three conservation equations hold independently.
There is no requirement that the three orientations agree on a common edge.

\paragraph{An exact parity test for the stronger method.}
For a partition $A_1,A_2,A_3$, all three complements are Eulerian exactly
when, at every vertex,
\[
       d_{A_i}(v)\equiv d_G(v)\pmod2\qquad(i=1,2,3).
\]
If $G$ is cubic, these three nonnegative degrees sum to three and must
all be odd. Hence each equals one. Thus in cubic graphs a partition
with Eulerian complements is equivalent to a proper three-edge-colouring.
This equivalence is proved directly; it is stronger than the target's
requirement that the complements support values of magnitudes up to three.

\paragraph{Checking the Petersen obstruction explicitly.}
Use vertices $u_i,v_i$ with $i\in\mathbb Z/5\mathbb Z$, outer edges
$u_iu_{i+1}$, spokes $u_iv_i$, and inner edges $v_iv_{i+2}$.
In a perfect matching the number of spokes is odd, because after removing
their endpoints both remaining five-cycles must have an even number of
vertices. It can therefore be one, three, or five. Three spokes leave
two indices that would have to be adjacent both in the outer cycle
(difference $\pm1$) and the inner cycle (difference $\pm2$), impossible.
With five spokes there is one matching. With one spoke, each remaining
four-vertex path has a unique perfect matching, giving five more.

For each of these six matchings the complement consists of two five-cycles.
For the all-spoke matching this is immediate. For a single-spoke matching,
rotation reduces to the spoke $u_0v_0$; direct tracing of the remaining
edges verifies the same lengths. The root batch script enumerates all
five-edge subsets and verifies both the count six and these cycle lengths.
In a proper three-edge-colouring, deleting any colour would leave cycles
alternating between the other two colours, so all would be even. Hence
the Petersen graph has no such colouring and no partition with all three
complements Eulerian.

\paragraph{What this does and does not rule out.}
The Petersen example rules out extending the unit-flow construction to
all bridgeless cubic graphs. A complement with a four-flow need not be
Eulerian: at a degree-three vertex values one and two can balance value
three. Consequently the parity obstruction does not apply to the target
statement. Allowing those non-unit values is precisely the flexibility
an improved argument must exploit, while avoiding bridges in each
complement and maintaining a single edge partition.

\paragraph{Final status.}
\textbf{UNRESOLVED.} The requested partition is constructed for all
graphs already admitting a four-flow and, explicitly, for every
three-edge-colourable cubic graph. The Petersen calculation identifies
an exact obstruction to the stronger Eulerian-complement strategy.
It supplies no counterexample to three four-flows and no general
construction for cubic graphs without three edge colours.

\subsection{Sources}
\begin{itemize}
\item[S1] ULAM.AI and the UnsolvedMath Contributors, supplied problems.json, record 3219 (OPG-132), OpenGarden collection. Catalogue identity and source transcription; CC BY 4.0.
\url{https://huggingface.co/datasets/ulamai/UnsolvedMath}.
\item[S2] Open Problem Garden, The three 4-flows conjecture. Original problem and discussion; the mathematical formulation below is expanded from the supplied source record.
\url{http://www.openproblemgarden.org/op/three_4_flows_conjecture}.

\end{itemize}
\end{document}
